% =====================================================================
%  CUERPO DEL MANUAL DEL SUPERVISOR - PharmaWeigh
% =====================================================================

\seccion{1. Acceso al sistema}

\tablaAcceso

\vspace{8pt}

Ingresa en \cd{https://pharma-formweigh.vercel.app} con tus credenciales.
Como supervisor tienes acceso a todos los módulos del operario, más las funciones
de aprobación de lotes, creación de recetas y firma electrónica.

% =====================================================================
\seccion{2. Gestión de inventario}

\subseccion{2.1 Aprobar o rechazar lotes}

Cuando llega materia prima a la planta, se registra con estado \cuarentena.
Es tu responsabilidad aprobar o rechazar los lotes después de la revisión de calidad.

\begin{enumerate}
    \item Ve a \menu{Inventario} $\rightarrow$ pestaña \menu{Lotes}.
    \item Identifica los lotes con estado \cuarentena.
    \item Revisa el certificado de análisis del proveedor.
    \item Haz clic en \boton{Aprobar} o \boton{Rechazar} según corresponda.
\end{enumerate}

\begin{cajaOjo}
\textbf{Un lote rechazado o en cuarentena NO puede usarse para dispensado.}
El sistema bloqueará cualquier intento de escanear un lote no aprobado durante
el proceso de pesaje. Esta validación es automática y no puede saltarse.
\end{cajaOjo}

\begin{cajaInfo}[Registro automático]
Cada cambio de estado de un lote se registra en la auditoría con tu nombre,
la fecha/hora, el estado anterior y el nuevo estado. Este registro es inmutable.
\end{cajaInfo}

\subseccion{2.2 Recepción de material nuevo}

\begin{enumerate}
    \item Ve a \menu{Inventario} $\rightarrow$ pestaña \menu{Recepción}.
    \item Escanea el \textbf{código de barras del material} con la pistola.
    \item El sistema muestra el nombre del material encontrado.
    \item Llena los datos del nuevo lote:
        \begin{itemize}
            \item \campo{Número de Lote} --- el identificador del proveedor
            \item \campo{Cantidad} --- peso o volumen recibido
            \item \campo{Proveedor} --- nombre del proveedor
            \item \campo{Fecha de Caducidad}
            \item \campo{Certificado de Análisis} (referencia)
        \end{itemize}
    \item Haz clic en \boton{Registrar Lote (Cuarentena)}.
    \item El lote se crea automáticamente en estado \cuarentena.
\end{enumerate}

\subseccion{2.3 Alertas de inventario}

Revisa regularmente:
\begin{itemize}
    \item \textbf{Materiales con stock bajo:} aparecen resaltados en rojo en la tabla de materiales.
    \item \textbf{Lotes próximos a caducar:} el dashboard muestra cuántos lotes caducan en 30 días.
    \item \textbf{Lotes en cuarentena:} no olvides aprobarlos para que estén disponibles.
\end{itemize}

% =====================================================================
\seccion{3. Gestión de recetas}

Como supervisor puedes crear y editar recetas (formulaciones).

\subseccion{3.1 Crear una receta nueva}

\begin{enumerate}
    \item Ve a \menu{Recetas} y haz clic en \boton{+ Nueva Receta}.
    \item Llena los datos generales:
        \begin{itemize}
            \item \campo{Código}: identificador único (ej: \codigo{REC-PCT-500})
            \item \campo{Nombre}: nombre descriptivo
            \item \campo{Rendimiento}: cantidad que produce la receta
            \item \campo{Unidad de Rendimiento}: kg, tabletas, L, etc.
        \end{itemize}
    \item Agrega ingredientes con \boton{+ Agregar Ingrediente}:
        \begin{itemize}
            \item \campo{Material}: selecciona de la lista de materiales registrados
            \item \campo{Cantidad}: peso/volumen exacto requerido
            \item \campo{Tolerancia Mín/Máx \%}: rango aceptable (ej: $-2\%$ / $+2\%$)
            \item \campo{Instrucciones}: texto que verá el operario en pantalla
            \item \campo{Material Peligroso}: activa la alerta roja en dispensado
        \end{itemize}
    \item El \textbf{orden de los ingredientes} define el orden de dispensado.
    \item Haz clic en \boton{Guardar Receta}.
\end{enumerate}

\begin{cajaTip}[Tolerancias]
Las tolerancias se expresan en porcentaje del target:
\begin{itemize}
    \item $\pm 1\%$ para principios activos (alta precisión)
    \item $\pm 2\%$ a $\pm 3\%$ para excipientes
    \item $\pm 5\%$ a $\pm 10\%$ para componentes menores (lubricantes, deslizantes)
\end{itemize}
El semáforo \semarojo{} bloquea la confirmación si el peso cae fuera de estos rangos.
\end{cajaTip}

% =====================================================================
\seccion{4. Órdenes de producción}

\subseccion{4.1 Crear una orden}

\begin{enumerate}
    \item Ve a \menu{Órdenes} y haz clic en \boton{+ Nueva Orden}.
    \item Selecciona la \campo{Receta} a producir.
    \item Define el \campo{Lote del Producto Final} (ej: \codigo{PROD-PCT-2024-003}).
    \item Define la \campo{Cantidad} (multiplicador sobre la receta base).
    \item Define la \campo{Prioridad}: Normal, Alta o Urgente.
    \item Haz clic en \boton{Crear Orden}.
\end{enumerate}

\subseccion{4.2 Estados de una orden}

\begin{center}
\begin{tabularx}{0.85\linewidth}{@{}l >{\raggedright\arraybackslash}X@{}}
\toprule
\textbf{Estado} & \textbf{Significado} \\
\midrule
\pendiente   & Creada, esperando que un operario inicie el dispensado. \\
\enproceso   & El operario está dispensando los ingredientes. \\
\dispensado  & Todos los ingredientes pesados y firmados por supervisor. \\
\completada  & Producción terminada. \\
\bottomrule
\end{tabularx}
\end{center}

% =====================================================================
\seccion{5. Firma electrónica de dispensados}

Esta es una de tus responsabilidades más críticas. Cuando un operario termina
de dispensar todos los ingredientes de una orden, el sistema te solicita tu firma.

\begin{enumerate}
    \item El operario te llama a su estación cuando aparece la pantalla de firma.
    \item \textbf{Verifica visualmente} que los contenedores correspondan.
    \item Revisa el \textbf{Resumen de pasos} en pantalla: verifica que todos los
          pesos estén dentro de tolerancia (palomitas verdes).
    \item Ingresa tu \campo{email} y tu \campo{contraseña}.
    \item Haz clic en \boton{Firmar y Aprobar}.
\end{enumerate}

\begin{cajaOjo}
\textbf{Tu firma es legalmente vinculante.} Al firmar certificas que:
\begin{itemize}
    \item Los materiales dispensados son correctos.
    \item Los pesos están dentro de las tolerancias aprobadas.
    \item El proceso se ejecutó siguiendo las BPM (Buenas Prácticas de Manufactura).
\end{itemize}
\textbf{NO firmes sin verificar.}
\end{cajaOjo}

% =====================================================================
\seccion{6. Auditoría}

Desde \menu{Auditoría} puedes consultar el registro completo de todas las
acciones del sistema.

\subseccion{6.1 Filtros disponibles}
\begin{itemize}
    \item \campo{Por acción}: LOGIN, CREAR\_RECETA, DISPENSAR, FIRMAR\_DISPENSADO, etc.
    \item \campo{Por usuario}: filtrar acciones de un operario específico.
\end{itemize}

\subseccion{6.2 Información registrada}
Cada entrada de auditoría contiene:
\begin{itemize}
    \item \textbf{Fecha y hora exacta} del evento.
    \item \textbf{Usuario} que ejecutó la acción y su rol.
    \item \textbf{Tipo de acción} (color codificado).
    \item \textbf{Entidad afectada} (tabla de la base de datos).
    \item \textbf{Detalles} del cambio (valores anteriores y nuevos).
\end{itemize}

\begin{cajaInfo}[21 CFR Part 11 --- Auditoría]
El log de auditoría es \textbf{inmutable}. No puede ser borrado, editado ni
desactivado. Esto cumple con los requisitos de la FDA para registros electrónicos.
En caso de inspección, este log es la evidencia principal de trazabilidad.
\end{cajaInfo}

% =====================================================================
\seccion{7. Gestión de usuarios}

Desde \menu{Usuarios} puedes consultar los usuarios del sistema.
Solo el \textbf{ADMIN} puede crear nuevos usuarios.

Si necesitas un nuevo operario o cambiar una contraseña, contacta al administrador.

% =====================================================================
\seccion{8. Preguntas frecuentes}

\subseccion{¿Puedo modificar una receta que ya tiene órdenes?}
Sí, pero se recomienda crear una nueva versión. Las órdenes existentes
mantienen la receta con la que fueron creadas.

\subseccion{¿Qué pasa si un operario dispensa fuera de tolerancia?}
El sistema lo bloquea automáticamente. El botón ``Confirmar Peso'' queda
deshabilitado y aparece un aviso rojo.

\subseccion{¿Puedo ver quién dispensó cada lote?}
Sí. En \menu{Auditoría} filtra por acción ``DISPENSAR''. Cada registro
muestra el operario, el material, el peso real vs.\ target, y si estuvo
dentro de tolerancia.

\subseccion{¿Cómo agrego un material nuevo?}
Ve a \menu{Inventario} $\rightarrow$ pestaña \menu{Materiales} $\rightarrow$
\boton{+ Nuevo Material}. Define el código de barras, nombre, unidad y stock mínimo.
